\documentclass[10pt,letterpaper,twocolumn]{article}

\usepackage[utf8]{inputenc}
\usepackage[T1]{fontenc}
\usepackage{newtxtext,newtxmath}
\usepackage[scaled=0.92]{helvet}
\usepackage[top=0.76in,bottom=0.82in,left=0.72in,right=0.72in,columnsep=0.23in]{geometry}
\usepackage{microtype}
\usepackage{amsmath,mathtools}
\usepackage{graphicx}
\usepackage{booktabs}
\usepackage{longtable}
\usepackage{array}
\usepackage{enumitem}
\usepackage{titlesec}
\usepackage{xcolor}
\usepackage[hidelinks]{hyperref}
\usepackage{fancyhdr}
\usepackage{setspace}

\setstretch{0.995}
\setlength{\parindent}{0.9em}
\setlength{\parskip}{0pt}
\setlength{\columnsep}{0.23in}
\setlist[itemize]{leftmargin=1.2em,itemsep=0.2em,topsep=0.35em}
\setlist[enumerate]{leftmargin=1.3em,itemsep=0.2em,topsep=0.35em}
\setlength{\abovedisplayskip}{5pt plus 2pt minus 2pt}
\setlength{\belowdisplayskip}{5pt plus 2pt minus 2pt}
\setlength{\abovedisplayshortskip}{4pt plus 2pt minus 2pt}
\setlength{\belowdisplayshortskip}{4pt plus 2pt minus 2pt}
\setlength{\textfloatsep}{8pt plus 2pt minus 2pt}
\setlength{\floatsep}{7pt plus 2pt minus 2pt}
\setlength{\intextsep}{7pt plus 2pt minus 2pt}
\setlength{\emergencystretch}{1.5em}
\raggedbottom

\definecolor{TitleRule}{HTML}{D9DEE7}
\definecolor{TitleNote}{HTML}{4B5563}
\definecolor{HeadingInk}{HTML}{1F2937}
\definecolor{HeadingRule}{HTML}{CBD5E1}

\renewcommand{\thesection}{\arabic{section}}
\renewcommand{\thesubsection}{\thesection.\arabic{subsection}}
\renewcommand{\thesubsubsection}{\thesubsection.\arabic{subsubsection}}
\setcounter{secnumdepth}{2}

\titleformat{\section}
  {\large\sffamily\bfseries\color{HeadingInk}}
  {\thesection}
  {0.65em}
  {}
\titleformat{name=\section,numberless}
  {\large\sffamily\bfseries\color{HeadingInk}}
  {}
  {0pt}
  {}
\titleformat{\subsection}
  {\normalsize\sffamily\bfseries\color{HeadingInk}}
  {\thesubsection}
  {0.55em}
  {}
\titleformat{\subsubsection}
  {\normalsize\sffamily\bfseries\color{HeadingInk}}
  {\thesubsubsection}
  {0.5em}
  {}
\titlespacing*{\section}{0pt}{0.95em}{0.35em}
\titlespacing*{\subsection}{0pt}{0.75em}{0.25em}
\titlespacing*{\subsubsection}{0pt}{0.55em}{0.2em}

\pagestyle{fancy}
\fancyhf{}
\fancyfoot[C]{\small\thepage}
\renewcommand{\headrulewidth}{0pt}

\providecommand{\tightlist}{%
  \setlength{\itemsep}{0pt}\setlength{\parskip}{0pt}}

\title{A Fourier-Domain Analysis of BHEX for Photon-Ring Inference}
\author{Jonathan R. Landers}
\date{}

\newcommand{\affiliationnote}{Independent researcher; not affiliated
with the BHEX collaboration or mission team.}

\makeatletter
\renewcommand{\maketitle}{%
  \twocolumn[{%
    \begin{@twocolumnfalse}
      \vspace*{-1.3em}
      {\LARGE\sffamily\bfseries\raggedright \@title \par}
      \vspace{0.35em}
      {\normalsize\sffamily\raggedright \@author \par}
      \vspace{0.16em}
      {\small\color{TitleNote}\raggedright \affiliationnote \par}
      \vspace{0.65em}
      {\color{TitleRule}\hrule height 0.8pt}
      \vspace{0.8em}
    \end{@twocolumnfalse}
  }]
}
\makeatother

\begin{document}
\maketitle

In recent work proposing the \textbf{Black Hole Explorer (BHEX)}
mission, the BHEX authors argue that sufficiently long-baseline
interferometric measurements can reveal, and in favorable regimes help isolate, the universal photon-ring
signature of a black hole from the more complicated emission of the
surrounding plasma [\hyperlink{ref-1}{1}, \hyperlink{ref-2}{2}]. In
practice, the measured visibility data contain both contributions, and
the observational challenge is to determine whether a structured
photon-ring signal can be reliably extracted from the mixture using
long-baseline amplitude signatures or more explicit structured
source-separation methods. This note develops one possible mathematical extension of
that problem framing, inspired by the BHEX visibility program.

The natural story is two-stage. First, one can ask how far a simple
long-baseline ring-oriented heuristic can be pushed before structured
plasma leakage creates bias. Second, one can ask whether that same
regime is still recoverable once the inverse problem is posed more
explicitly. The first question leads to a mismatch bound between a
direct amplitude-readout heuristic and a plasma-aware estimator. The
second leads to a recoverability statement: the direct amplitude readout
can lose clarity before ring non-recoverability is reached.

\hypertarget{a-first-comparison-template-mismatch-versus-plasma-aware-projection}{%
\section{A first comparison: template mismatch versus plasma-aware
projection}\label{a-first-comparison-template-mismatch-versus-plasma-aware-projection}}

\par\smallskip\noindent{}\textbf{Lemma (Mismatch between heuristic ring
detection and structured plasma-aware separation).} Let
\(y \in \mathcal H\) be observed visibility data in a complex Hilbert
space, decomposed as

\[
y = \alpha g + p + \varepsilon,
\]

where \(g \in \mathcal H\) is a photon-ring template,
\(\alpha \in \mathbb C\) is the ring amplitude,
\(p \in \mathcal P \subset \mathcal H\) is a plasma component lying in a
plasma subspace \(\mathcal P\), and \(\varepsilon \in \mathcal H\) is
noise. Define the heuristic estimator

\[
\hat\alpha_{\mathrm{heur}} = \frac{\langle y,g\rangle}{\|g\|^2},
\]

and the plasma-aware estimator

\[
\hat\alpha_{\mathrm{model}}
=
\frac{\left\langle (I-\Pi_{\mathcal P})y,\; (I-\Pi_{\mathcal P})g\right\rangle}
{\|(I-\Pi_{\mathcal P})g\|^2},
\]

where \(\Pi_{\mathcal P}\) denotes the orthogonal projector onto
\(\mathcal P\). Then

\[
\bigl|\hat\alpha_{\mathrm{heur}}-\hat\alpha_{\mathrm{model}}\bigr|
\le
\mu(g,\mathcal P)\frac{\|p\|}{\|g\|}
+
\frac{\|\varepsilon\|}{\|g\|}
+
\frac{\|(I-\Pi_{\mathcal P})\varepsilon\|}{\|(I-\Pi_{\mathcal P})g\|},
\]

where

\[
\mu(g,\mathcal P)=\frac{\|\Pi_{\mathcal P}g\|}{\|g\|}
\]

is the ring-plasma coherence.

\par\smallskip\noindent{}\textbf{Proof.} First, by direct substitution,

\[
\hat\alpha_{\mathrm{heur}}
=
\frac{\langle \alpha g + p + \varepsilon,\; g\rangle}{\|g\|^2}
=
\alpha
+
\frac{\langle p,g\rangle}{\|g\|^2}
+
\frac{\langle \varepsilon,g\rangle}{\|g\|^2}.
\]

Hence

\[
\hat\alpha_{\mathrm{heur}}-\alpha
=
\frac{\langle p,g\rangle}{\|g\|^2}
+
\frac{\langle \varepsilon,g\rangle}{\|g\|^2}.
\]

Now define

\[
g_\perp = (I-\Pi_{\mathcal P})g.
\]

Since \(p \in \mathcal P\), we have

\[
(I-\Pi_{\mathcal P})p = 0.
\]

Therefore

\[
(I-\Pi_{\mathcal P})y
=
\alpha (I-\Pi_{\mathcal P})g + (I-\Pi_{\mathcal P})\varepsilon
=
\alpha g_\perp + (I-\Pi_{\mathcal P})\varepsilon.
\]

Substituting into the plasma-aware estimator gives

\[
\hat\alpha_{\mathrm{model}}
=
\frac{\langle \alpha g_\perp + (I-\Pi_{\mathcal P})\varepsilon,\; g_\perp\rangle}{\|g_\perp\|^2}
=
\alpha
+
\frac{\langle (I-\Pi_{\mathcal P})\varepsilon,\; g_\perp\rangle}{\|g_\perp\|^2}.
\]

So

\[
\hat\alpha_{\mathrm{model}}-\alpha
=
\frac{\langle (I-\Pi_{\mathcal P})\varepsilon,\; g_\perp\rangle}{\|g_\perp\|^2}.
\]

Subtracting the two estimators yields

\[
\hat\alpha_{\mathrm{heur}}-\hat\alpha_{\mathrm{model}}
=
\frac{\langle p,g\rangle}{\|g\|^2}
+
\frac{\langle \varepsilon,g\rangle}{\|g\|^2}
-
\frac{\langle (I-\Pi_{\mathcal P})\varepsilon,\; g_\perp\rangle}{\|g_\perp\|^2}.
\]

Applying the triangle inequality,

\[
\bigl|\hat\alpha_{\mathrm{heur}}-\hat\alpha_{\mathrm{model}}\bigr|
\le
\frac{|\langle p,g\rangle|}{\|g\|^2}
+
\frac{|\langle \varepsilon,g\rangle|}{\|g\|^2}
+
\frac{|\langle (I-\Pi_{\mathcal P})\varepsilon,\; g_\perp\rangle|}{\|g_\perp\|^2}.
\]

Using Cauchy--Schwarz on the last two terms,

\[
\frac{|\langle \varepsilon,g\rangle|}{\|g\|^2}
\le
\frac{\|\varepsilon\|\,\|g\|}{\|g\|^2}
=
\frac{\|\varepsilon\|}{\|g\|},
\]

and

\[
\frac{|\langle (I-\Pi_{\mathcal P})\varepsilon,\; g_\perp\rangle|}{\|g_\perp\|^2}
\le
\frac{\|(I-\Pi_{\mathcal P})\varepsilon\|\,\|g_\perp\|}{\|g_\perp\|^2}
=
\frac{\|(I-\Pi_{\mathcal P})\varepsilon\|}{\|g_\perp\|}.
\]

It remains to bound the plasma term. Since \(p \in \mathcal P\),

\[
\langle p,g\rangle = \langle p,\Pi_{\mathcal P}g\rangle,
\]

so again by Cauchy--Schwarz,

\[
|\langle p,g\rangle|
=
|\langle p,\Pi_{\mathcal P}g\rangle|
\le
\|p\|\,\|\Pi_{\mathcal P}g\|.
\]

Thus

\[
\frac{|\langle p,g\rangle|}{\|g\|^2}
\le
\frac{\|p\|\,\|\Pi_{\mathcal P}g\|}{\|g\|^2}
=
\left(\frac{\|\Pi_{\mathcal P}g\|}{\|g\|}\right)\frac{\|p\|}{\|g\|}
=
\mu(g,\mathcal P)\frac{\|p\|}{\|g\|}.
\]

Combining the three bounds gives

\[
\bigl|\hat\alpha_{\mathrm{heur}}-\hat\alpha_{\mathrm{model}}\bigr|
\le
\mu(g,\mathcal P)\frac{\|p\|}{\|g\|}
+
\frac{\|\varepsilon\|}{\|g\|}
+
\frac{\|(I-\Pi_{\mathcal P})\varepsilon\|}{\|(I-\Pi_{\mathcal P})g\|}.
\]

This proves the claim. \(\square\)

\par\smallskip\noindent{}\textbf{Interpretation.} This lemma makes
precise the gap between a template-based long-baseline heuristic and a
plasma-aware structured estimator in the same general spirit as the BHEX
visibility-analysis program [\hyperlink{ref-2}{2}]. The mismatch is
controlled by three interpretable quantities: ring-plasma coherence,
plasma magnitude, and measurement noise. The more the photon-ring
template overlaps realistic plasma structure, the more the heuristic
estimator can be biased relative to a structured projection.

\hypertarget{from-mismatch-to-recoverability}{%
\section{From mismatch to
recoverability}\label{from-mismatch-to-recoverability}}

The previous lemma quantifies the discrepancy between two estimators,
but it does not yet answer the more interesting question: if the direct
amplitude readout becomes hard to interpret, does recoverability
necessarily disappear with it? The next proposition formalizes the
stricter and more useful claim that the answer can be no. In other
words, the heuristic readability limit can arrive before the
source-separation limit.

\par\smallskip\noindent{}\textbf{Proposition (direct-readout limit
before separation limit).} Let the observed visibility be

\[
y=\alpha g_\theta + q + \varepsilon,
\]

in a Hilbert space \(\mathcal H\), where:

\begin{itemize}
\tightlist
\item
  \(g_\theta \in \mathcal G\) is the photon-ring visibility family,
  parameterized by ring shape \(\theta\),
\item
  \(\alpha \in \mathbb C\) is the ring amplitude,
\item
  \(q \in \mathcal Q\) is a structured nuisance term representing
  plasma, finite-width \(n=1\) distortion, scattering, or averaging
  artifacts,
\item
  \(\varepsilon\) is noise.
\end{itemize}

Assume, as a simplified abstraction motivated by the BHEX visibility
program, that a long-baseline amplitude heuristic estimates \(\theta\) from
the oscillatory structure of the \textbf{visibility amplitude} \(|y|\)
[\hyperlink{ref-2}{2}]. Assume also that a structured estimator is given
by

\[
(\hat\alpha,\hat\theta,\hat q)
\in
\arg\min_{\alpha,\theta,q\in\mathcal Q}
\|y-\alpha g_\theta-q\|_{\mathcal H}^2+\lambda R(q),
\]

where \(R\) encodes nuisance structure. Finally, assume a local
identifiability condition: for some \(c>0\),

\[
\|\alpha g_\theta + q - \alpha' g_{\theta'} - q'\|_{\mathcal H}
\ge
c\Big(|\alpha-\alpha'| + d(\theta,\theta') + \|q-q'\|_{\mathcal H}\Big)
\]

for all nearby triples \((\alpha,\theta,q)\) and
\((\alpha',\theta',q')\).

Then there exists a nonempty regime in which the amplitude-based
heuristic becomes biased or visually unstable, but the structured
estimator still recovers \(\theta\) stably.

\par\smallskip\noindent{}\textbf{Proof.} Write

\[
y=\alpha g_\theta + r,
\qquad r:=q+\varepsilon.
\]

The heuristic extracts \(\theta\) from extrema or oscillation spacing of
\(|y|\) on long baselines. It is therefore accurate only when

\[
|y|=|\alpha g_\theta+r|
\]

has nearly the same oscillatory pattern as \(|\alpha g_\theta|\). But
modulus is nonlinear, so even moderate perturbations can shift extrema
and spacing. Pointwise,

\[
\bigl||\alpha g_\theta+r|-|\alpha g_\theta|\bigr|
\le |r|,
\]

yet small amplitude perturbation does \textbf{not} imply small
perturbation of extrema locations. Near a shallow oscillation peak, an
\(O(\delta)\) perturbation can induce an \(O(\delta/\kappa)\)
peak-location error, where \(\kappa\) denotes local curvature. Thus the
heuristic can lose direct-readout accuracy once structured contamination
distorts the oscillatory pattern, even if the ring signal is still
present.

By contrast, the structured estimator fits the full model
\(y=\alpha g_\theta+q+\varepsilon\). Under the identifiability
inequality, standard stability of constrained least squares implies that
if \(\|\varepsilon\|_{\mathcal H}\le \eta\), then

\[
|\hat\alpha-\alpha| + d(\hat\theta,\theta) + \|\hat q-q\|_{\mathcal H}
\le \frac{C}{c}\,\eta
\]

for some constant \(C\), since any competing triple separated from the
truth by more than \(O(\eta)\) would violate the lower bound and
increase the residual.

Therefore there is a regime where:

\begin{itemize}
\tightlist
\item
  the nuisance \(q\) is large enough to distort the amplitude
  oscillation and make direct peak-reading unreliable,
\item
  but still small enough, and sufficiently structured, that the full
  inverse problem remains identifiable and stable.
\end{itemize}

Hence

\[
\begin{aligned}
\{\text{cleanly amplitude-readable}\}
&\subsetneq \{\text{structured-separation recoverable}\}
\end{aligned}
\]

is possible. \(\square\)

\par\smallskip\noindent{}\textbf{One-line interpretation.} The BHEX
heuristic is especially powerful when the ring is \textbf{cleanly
readable} in amplitude space [\hyperlink{ref-2}{2}]; the structured
estimator succeeds under the weaker condition that the ring remains
\textbf{identifiable} after joint modeling of structured contamination.

\hypertarget{physical-picture}{%
\section{Physical picture}\label{physical-picture}}

At the image level, the distinction is intuitive. The observed
black-hole image is typically dominated by a \textbf{broader plasma
emission ring or crescent}, produced by synchrotron emission from hot
plasma, while the \textbf{photon ring} is a much thinner structure at
nearly the same diameter, set more directly by strong lensing geometry
[\hyperlink{ref-2}{2}, \hyperlink{ref-3}{3}]. This means the source
naturally presents a broad, bright, plasma-dominated feature, while the
measurement goal is to isolate a much thinner, more universal ring
hidden within or beneath it. In Fourier space, that geometric difference
becomes a problem of structured overlap: broad plasma emission and thin
photon-ring oscillations can coexist on the same visibility domain, but
need not be equally identifiable.

This is exactly why the source-separation perspective is useful as an
extension alongside the BHEX program. The heuristic says: go where the
photon ring should show itself most directly. The structured method
says: even if that regime is not perfectly clean, the ring may still be
recoverable provided its structure remains distinguishable from the
nuisance family.

\hypertarget{connection-to-noise}{%
\section{Connection to noise}\label{connection-to-noise}}

The first lemma separates the mismatch into a \textbf{plasma leakage
term} and two \textbf{noise terms}. The first noise term,
\(\|\varepsilon\|/\|g\|\), is the raw sensitivity of the heuristic
estimator to measurement noise, while the second,

\[
\frac{\|(I-\Pi_{\mathcal P})\varepsilon\|}{\|(I-\Pi_{\mathcal P})g\|},
\]

shows that the structured method is governed by the component of the
noise that survives plasma projection, normalized by the strength of the
projected ring template. Thus the plasma-aware method reduces bias from
source confusion, but its advantage is greatest when the projected ring
signal remains sufficiently strong relative to the projected noise.

The proposition sharpens this picture. Noise and nuisance structure do
not merely degrade the signal-to-noise ratio; they can deform the
oscillatory pattern used by the heuristic before they destroy
identifiability altogether. This is the conceptual gap between a direct
long-baseline readout strategy and a more robust inverse-problem
formulation.

\hypertarget{astrophysical-relevance}{%
\section{Astrophysical relevance}\label{astrophysical-relevance}}

The mathematical content becomes most compelling when the abstract
nuisance class is tied to realistic plasma and propagation effects. A
natural next step is therefore to estimate either the ring-plasma
coherence \(\mu(g,\mathcal P)\) or a more general identifiability margin
for nuisance classes informed by GRMHD simulations, radiative transfer
calculations, scattering models, or empirical variability constraints
[\hyperlink{ref-2}{2}, \hyperlink{ref-3}{3}]. Such a result would
connect the abstract bounds directly to the astrophysical conditions
under which long-baseline photon-ring detection is expected to succeed,
become biased, or remain recoverable only through structured separation.

In that sense, the BHEX heuristic can be viewed as an elegant front-end
to the inference problem: it targets the regime where the photon ring
more or less announces itself [\hyperlink{ref-2}{2}]. The
source-separation formalism is a complementary robustness layer: it
targets the broader regime in which the data are less cleanly readable,
but the ring still leaves enough structured trace to be recovered.

\hypertarget{key-interesting-points}{%
\section{Key interesting points}\label{key-interesting-points}}

\begin{itemize}
\tightlist
\item
  The writeup develops a
  \textbf{structured source-separation abstraction in Fourier space}
  inspired by BHEX photon-ring inference [\hyperlink{ref-1}{1}, \hyperlink{ref-2}{2}].
\item
  The first lemma gives a \textbf{clean mismatch bound} between a direct
  amplitude-readout heuristic and a plasma-aware estimator, controlled
  by ring-plasma coherence, plasma magnitude, and noise.
\item
  The new proposition makes the stronger point that the
  \textbf{heuristic readability limit can occur before the
  recoverability limit} in this abstraction, provided the ring remains identifiable under
  joint modeling.
\item
  The physical picture is that the image naturally shows a \textbf{broad
  plasma ring or crescent}, while the scientific target is a
  \textbf{much thinner photon ring} at nearly the same diameter
  [\hyperlink{ref-2}{2}, \hyperlink{ref-3}{3}].
\item
  The most promising bridge to astrophysics is to estimate coherence or
  identifiability margins for \textbf{GRMHD-informed nuisance classes},
  linking this abstract theory to real mission regimes anticipated by
  BHEX forecast work [\hyperlink{ref-2}{2}, \hyperlink{ref-3}{3}].
\end{itemize}

\hypertarget{references}{%
\section*{References}\label{references}}
\addcontentsline{toc}{section}{References}

\noindent{}\hypertarget{ref-1}{}{[}1{]} Black Hole Explorer
Collaboration, \emph{Black Hole Explorer: Motivation and Vision},
arXiv:2406.12917. https://arxiv.org/abs/2406.12917

\noindent{}\hypertarget{ref-2}{}{[}2{]} Black Hole Explorer
Collaboration, \emph{The Black Hole Explorer: Photon Ring Science, Detection and Shape Measurement}, arXiv:2406.09498. https://arxiv.org/abs/2406.09498

\noindent{}\hypertarget{ref-3}{}{[}3{]} Black Hole Explorer
Collaboration, \emph{Interferometric Inference of Black Hole Spin from
Photon Ring Size and Brightness}, arXiv:2509.23628.
https://arxiv.org/abs/2509.23628

\end{document}
